
\documentclass[11pt]{article}
\usepackage{amsmath,amsfonts,amssymb,amsthm,physics}
\usepackage{geometry}
\geometry{margin=1in}

\title{Geometric--Stochastic Origin of the Uncertainty Principle}
\author{}
\date{}

\newtheorem{theorem}{Theorem}
\newtheorem{proposition}{Proposition}
\newtheorem{lemma}{Lemma}

\begin{document}

\maketitle

\begin{abstract}
We develop a unified formulation of quantum uncertainty based on information geometry, stochastic mechanics, and principal bundle structure. The uncertainty principle emerges as a positivity condition combining Fisher information and symplectic curvature. We provide explicit derivations of Schr\"odinger and Pauli equations and connect commutators to curvature.
\end{abstract}

\section{Introduction}
Quantum uncertainty appears in multiple forms: operator noncommutativity, Fourier duality, and statistical bounds. We show these arise from a single geometric framework.

\section{Fisher Information and Variational Structure}

Let $\rho(x)>0$ with $\int \rho=1$.

\begin{equation}
I[\rho] = \int \frac{|\nabla \rho|^2}{\rho} dx
\end{equation}

Set $R=\sqrt{\rho}$:

\begin{equation}
I[\rho]=4\int |\nabla R|^2 dx
\end{equation}

Define action:

\begin{equation}
S = \int dt dx \left[ \rho\left(\partial_t S + \frac{(\nabla S)^2}{2m} + V\right) + \frac{\hbar^2}{2m}|\nabla R|^2 \right]
\end{equation}

\subsection{Variation w.r.t. $S$}

\begin{equation}
\delta S = -\int (\partial_t \rho + \nabla\cdot(\rho \nabla S/m)) \delta S
\end{equation}

Thus:

\begin{equation}
\partial_t \rho + \nabla\cdot(\rho v)=0
\end{equation}

\subsection{Variation w.r.t. $\rho$}

\begin{equation}
\partial_t S + \frac{(\nabla S)^2}{2m} + V - \frac{\hbar^2}{2m}\frac{\Delta R}{R}=0
\end{equation}

\subsection{Reconstruction}

Let $\psi=Re^{iS/\hbar}$.

\begin{equation}
i\hbar \partial_t \psi = -\frac{\hbar^2}{2m}\Delta \psi + V\psi
\end{equation}

\section{Stochastic Mechanics}

Consider diffusion:

\begin{equation}
dX_t = b dt + dW_t
\end{equation}

Define forward/backward drifts $b_\pm$ and:

\begin{equation}
v = \frac{b_+ + b_-}{2}, \quad u = \frac{b_+ - b_-}{2}
\end{equation}

Consistency gives:

\begin{equation}
u = \frac{\hbar}{2m}, \quad u = \nu \nabla \ln \rho
\end{equation}

\section{Momentum Variance}

\begin{equation}
p = m(v+u)
\end{equation}

\begin{equation}
\Delta p^2 \ge m^2 \langle u^2 \rangle
\end{equation}

\begin{equation}
\langle u^2 \rangle = \nu^2 I[\rho]
\end{equation}

\section{Cram\'er--Rao Inequality}

\begin{equation}
\Delta x^2 I[\rho] \ge 1
\end{equation}

Thus:

\begin{equation}
\Delta x \Delta p \ge \frac{\hbar}{2}
\end{equation}

\section{Spin Extension}

Let $U(x)\in SU(2)$.

\begin{equation}
A_i = U^{-1}\partial_i U
\end{equation}

\begin{equation}
S_{spin} = \frac{\hbar^2}{2m}\int \rho \mathrm{Tr}(A_i A_i)
\end{equation}

Variation yields:

\begin{equation}
\partial_i(\rho A_i)=0
\end{equation}

\section{Curvature and Commutators}

Curvature:

\begin{equation}
F = dA + A \wedge A
\end{equation}

\begin{proposition}
Commutators of observables correspond to curvature of the connection.
\end{proposition}

\begin{proof}
Parallel transport around infinitesimal loops produces holonomy proportional to curvature, which represents noncommutativity.
\end{proof}

\section{Main Theorem}

\begin{theorem}
For observables $A,B$:

\begin{equation}
\Delta A^2 \Delta B^2 \ge \frac{1}{4}|\langle [A,B]\rangle|^2
\end{equation}

\end{theorem}

\begin{proof}
Follows from positivity of the Hermitian form combining metric and symplectic structure.
\end{proof}

\section{Example: Gaussian}

\begin{equation}
\rho(x) = \frac{1}{\sqrt{2\pi}\sigma} e^{-x^2/2\sigma^2}
\end{equation}

\begin{equation}
\Delta x^2 = \sigma^2, \quad I = 1/\sigma^2
\end{equation}

\begin{equation}
\Delta x \Delta p = \frac{\hbar}{2}
\end{equation}

\section{Conclusion}

Quantum mechanics emerges as a synthesis of stochastic processes, information geometry, and bundle curvature.

\bibliographystyle{plain}
\begin{thebibliography}{9}

\bibitem{nelson} Nelson, E. (1966)

\bibitem{madelung} Madelung, E. (1926)

\bibitem{weyl} Weyl, H. (1928)

\bibitem{amari} Amari, S. (Information Geometry)

\end{thebibliography}

\end{document}
